\documentclass[11pt]{amsart}
\usepackage[margin=1in]{geometry}
\usepackage{lmodern,microtype,amsmath,amssymb,amsthm,needspace}
\usepackage[hidelinks]{hyperref}
\newtheorem{theorem}{Theorem}[section]
\newtheorem{proposition}[theorem]{Proposition}
\newtheorem{lemma}[theorem]{Lemma}
\newtheorem{corollary}[theorem]{Corollary}
\theoremstyle{definition}
\newtheorem{definition}[theorem]{Definition}
\newtheorem{remark}[theorem]{Remark}
\newcommand{\R}{\mathbb R}
\DeclareMathOperator{\supp}{supp}
\DeclareMathOperator{\Vol}{Vol}
\emergencystretch=2em
\title{Isometric tight tori with different enclosed volumes}
\author{Yuetong Zhang}
\address{Peking University}
\email{2300010810@stu.pku.edu.cn}
\email{gzhangyuetong@gmail.com}
\date{October 29, 2026}
\hypersetup{pdftitle={Isometric tight tori with different enclosed volumes},pdfauthor={Yuetong Zhang}}
\begin{document}
\begin{abstract}
We construct two smooth embedded tight tori in Euclidean three-space
whose induced metrics agree and whose enclosed volumes differ. The
construction uses one closed asymptotic curve from the annulus of
Banks, Ghomi, Qiu, and Raffaelli. We modify the ends of that annulus
to infinite-order planar contact and complete it to a tight torus.
For any tight torus with one closed asymptotic curve in negative
curvature, we prove that a modification supported near the curve
produces an infinitesimal bending with nonzero normal flux.
Gauss-map injectivity forces an inflection, a local change of the curve
arranges a derivative condition, and a choice of rulings makes the
nearby asymptotic curves periodic. A bending on this ruled annulus is
then extended by zero to the completed surface. Opposite multiples of
the bending give the required pair. Each stage is stated as a separate
existence or implication theorem with explicit hypotheses.
\end{abstract}
\maketitle

\section{Definitions, main result, and the three implications}

All manifolds, maps, and fields are smooth, meaning $C^\infty$, unless
another regularity is specified. A \emph{closed surface} is a compact
connected smooth two-dimensional manifold without boundary. We use
the standard inner product, orientation, cross product, and Lebesgue
volume in $\R^3$. Write $T^2=(\R/2\pi\mathbb Z)^2$.
An \emph{embedding} is an injective immersion that is a homeomorphism
onto its image.

For an immersion $X:M\to\R^3$, its induced metric is
\[
 g_X(v,w)=\langle dX_p(v),dX_p(w)\rangle
 \quad (p\in M,\ v,w\in T_pM).
\]
Write $K_X$ and $dA_X$ for the Gaussian curvature and area measure of
this metric. If $n$ is a unit normal, our second fundamental form has
local coefficients
\[
 \mathrm{II}_{ij}=\langle X_{ij},n\rangle.
\]
The normal component makes this definition independent of the chosen
coordinates. In a coordinate chart, if $g_{ij}=\langle X_i,X_j\rangle$,
then
\begin{equation}\label{eq:gauss-curvature}
 K_X=\frac{\mathrm{II}_{11}\mathrm{II}_{22}
                 -\mathrm{II}_{12}^2}{g_{11}g_{22}-g_{12}^2}.
\end{equation}
This is the Gauss equation; see \cite[\S3-3 and \S4-3]{doCarmo}.
Reversing $n$ changes $\mathrm{II}$ by a sign and leaves $K_X$ unchanged.

\begin{definition}\label{def:tight}
An embedding $X:M\to\R^3$ of a closed surface is \emph{tight} if, for
every $e\in\R^3\setminus\{0\}$ and $t\in\R$, the set
\[
 X(M)\cap\{x\in\R^3:\langle x,e\rangle>t\}
\]
is empty or connected. This is the two-piece property for surfaces.

By the Jordan--Brouwer separation theorem, the image of a closed
embedded surface in $\R^3$ separates space into a bounded and an
unbounded component. Its \emph{enclosed volume}
$\Vol(X)$ is the Lebesgue volume of the bounded component, and
$\nu_X$ denotes its outward unit normal.

Two embeddings of the same source are \emph{isometric by the identity}
if their induced metrics are equal. Their images are \emph{congruent}
if a Euclidean isometry of $\R^3$ maps one image onto the other.
\end{definition}

\begin{definition}\label{def:bending}
A smooth map $Y:M\to\R^3$ is an \emph{infinitesimal bending} of an
immersion $X$ if
\begin{equation}\label{eq:strain}
 \langle dX_p(v),dY_p(w)\rangle+
 \langle dX_p(w),dY_p(v)\rangle=0
 \quad(p\in M,\ v,w\in T_pM).
\end{equation}
It is a vector field along $X$, and is not required to be tangent to
$X(M)$. Its \emph{support} is the closure in $M$ of
$\{p:Y(p)\ne0\}$. For a closed embedded surface its \emph{outward normal
flux} is
\[
 \int_M\langle Y,\nu_X\rangle\,dA_X.
\]
For subsets $A\subset U\subset M$, the notation $A\Subset U$ means that
$\overline A$ is compact and contained in the open set $U$.
The support of $X-X_0$ always refers to the pointwise difference of
two maps on the same source.
\end{definition}

\begin{definition}\label{def:asymptotic}
A regular curve $c$ on an immersed surface $X$ is \emph{asymptotic}
if $\mathrm{II}(c',c')=0$. Equivalently, for its image curve
$\gamma=X\circ c$, the acceleration $\gamma''$ is tangent to the
surface. For any change of regular parameter, the extra term in
the acceleration is a multiple of $\gamma'$, so this condition is
independent of the parameter. When $X$ is an embedding, a
\emph{closed embedded asymptotic curve} means the image of a smooth
embedding $c:S^1\to M$ with this property.
\end{definition}

\begin{theorem}[Main result]\label{thm:main}
There exist a smooth tight embedding $X:T^2\to\R^3$ and an
infinitesimal bending $Y:T^2\to\R^3$ such that
\[
 \supp Y\Subset\{K_X<0\},
 \qquad \int_{T^2}\langle Y,\nu_X\rangle\,dA_X\ne0.
\]
There is $a_0>0$ such that for every $0<a<a_0$ the maps
$X_\pm=X\pm aY$ are smooth tight embeddings satisfying
\begin{equation}\label{eq:main-pair}
 g_{X_+}=g_{X_-},\qquad \Vol(X_+)\ne\Vol(X_-).
\end{equation}
Consequently their images are noncongruent. Both maps agree with $X$
outside $\supp Y$, and in particular on a nonempty open set.
\end{theorem}

Here equality of metrics gives an intrinsic isometry by the identity
map of $T^2$. We do not assert that the metrics of $X+tY$ are independent
of $t$. The volume-stationarity problem for infinitesimal bendings is
discussed by V.~A.~Aleksandrov \cite{Aleksandrov}.

The proof is organized through the following three statements. Their
inputs and outputs are geometric objects on specified sources, so that
the existence construction and the local deformation can be used
independently.

\begin{theorem}[Existence of a starting torus]\label{thm:seed}
There exist a smooth tight embedding $X_0:T^2\to\R^3$ and a smooth
embedding $c:S^1\to T^2$ such that $X_0\circ c$ is asymptotic and
$K_{X_0}<0$ on $c(S^1)$.
\end{theorem}

Theorem~\ref{thm:seed} is obtained in Section~\ref{sec:seed} from the
construction of Banks--Ghomi--Qiu--Raffaelli \cite{BGQR}. Only one of
their closed asymptotic curves is used. The completion includes planar
regions, which are allowed in the present statement and remain fixed
under the subsequent deformation.

\begin{theorem}[One curve implies a supported bending]\label{thm:one-curve}
Suppose $X_0:T^2\to\R^3$ is a smooth tight embedding and
$c:S^1\to T^2$ is a smooth embedding such that $X_0\circ c$ is
asymptotic and $K_{X_0}<0$ along $c(S^1)$.
For every open set $U$ with
\[
 c(S^1)\subset U\subset\{K_{X_0}<0\},
\]
and every $\varepsilon>0$, there exist a smooth tight embedding
$X:T^2\to\R^3$ and an infinitesimal bending $Y$ of $X$ such that
\begin{equation}\label{eq:one-curve-output}
 \begin{gathered}
 \supp(X-X_0)\Subset U,\qquad K_X<0\ \hbox{on }U,\qquad
 \|X-X_0\|_{C^1}<\varepsilon,\\
 \supp Y\Subset U,\qquad
 \int_{T^2}\langle Y,\nu_X\rangle\,dA_X\ne0.
 \end{gathered}
\end{equation}
The $C^1$ norm is computed using any fixed smooth Riemannian metric
on the compact source and the Euclidean norm in the target.
The same convention, with covariant derivatives for the chosen
background metric, is used for $C^2$ norms below.
\end{theorem}

The bending in this theorem belongs to the modified embedding $X$.
No condition on derivatives of the torsion, and no family of closed
asymptotic curves, is required of $X_0\circ c$.
Sections~\ref{sec:frame}--\ref{sec:insertion} prove this theorem.

\begin{proposition}[A supported bending implies an isometric pair]
\label{prop:pair}
Suppose $X:M\to\R^3$ is a smooth tight embedding of a closed surface,
and $Y$ is an infinitesimal bending with
\[
 \supp Y\Subset\{K_X<0\},\qquad
 \int_M\langle Y,\nu_X\rangle\,dA_X\ne0.
\]
Then, for all sufficiently small $a>0$, $X_\pm=X\pm aY$ are tight
embeddings with unequal enclosed volumes and
\begin{equation}\label{eq:exact-metric}
 g_{X_\pm}=g_X+a^2\langle dY,dY\rangle,
 \quad
 \langle dY,dY\rangle(v,w)=\langle dY(v),dY(w)\rangle.
\end{equation}
\end{proposition}

Proposition~\ref{prop:pair} is proved in Section~\ref{sec:reductions}.
Theorem~\ref{thm:seed} supplies exactly the input of
Theorem~\ref{thm:one-curve}; its output supplies exactly the input of
Proposition~\ref{prop:pair}. This separates the global existence
problem from the local bending problem.

\begin{proof}[Proof of Theorem~\ref{thm:main} from these statements]
Choose $X_0,c$ by Theorem~\ref{thm:seed}. Compactness of $c(S^1)$
and continuity of $K_{X_0}$ give an open neighborhood $U$ with
$\overline U\subset\{K_{X_0}<0\}$. Apply
Theorem~\ref{thm:one-curve} and then Proposition~\ref{prop:pair}.
The positive-curvature locus of a tight surface is nonempty by
Theorem~\ref{thm:tightness} below, and is unchanged by these
modifications. It supplies the stated common open set.
\end{proof}

\section{Classical reductions and the volume argument}
\label{sec:reductions}

\begin{theorem}[Classical curvature criterion for tightness]
\label{thm:tightness}
For a smooth embedding $X:M\to\R^3$ of a closed surface, the two-piece
property of Definition~\ref{def:tight} is equivalent to
\begin{equation}\label{eq:tightness}
 \int_M\max(K_X,0)\,dA_X=4\pi.
\end{equation}
\end{theorem}

We use this classical theorem as an external result; see
Banchoff--K\"uhnel \cite[Definition~1.1.1, p.~57]{BK}.
By Gauss--Bonnet, its equivalent
total-absolute-curvature formulation is
$\int_M|K_X|\,dA_X=8\pi-2\pi\chi(M)$.
This criterion is useful here because it reduces preservation of
tightness to preservation of the positive-curvature part. It allows
zero-curvature regions.

\begin{lemma}[Preservation under a negative-curvature modification]
\label{lem:tight-preservation}
Suppose $X_0:M\to\R^3$ is a tight embedding of a closed surface,
$U\subset\{K_{X_0}<0\}$ is open, and $X:M\to\R^3$ is an embedding
with $\supp(X-X_0)\Subset U$ and $K_X<0$ on $U$.
Then $X$ is tight.
In particular these conditions hold for every sufficiently small
$C^2$ perturbation supported in a fixed compact subset of
$\{K_{X_0}<0\}$.
\end{lemma}
\begin{proof}
On $U$ both positive curvature densities vanish. Outside the compact
support of $X-X_0$, the maps agree on a neighborhood and hence have
the same metric and curvature. Thus their total positive curvatures
are equal, and Theorem~\ref{thm:tightness} applies.
For the last assertion choose a compact neighborhood of the prescribed
support contained in $\{K_{X_0}<0\}$. On this neighborhood the original
curvature has a strictly negative upper bound. Formula
\eqref{eq:gauss-curvature} and continuity in the first two derivatives
preserve that bound under a sufficiently small $C^2$ change.
Embeddings of a compact source form an open set in the $C^1$
topology \cite[Chapter~2, \S1, Theorem~1.4]{Hirsch}.
\end{proof}

\begin{lemma}[First variation of enclosed volume]\label{lem:volume}
For an embedding $X:M\to\R^3$ of a closed surface and a smooth field
$Y:M\to\R^3$, the maps $X+tY$ are embeddings for sufficiently small
$|t|$, and
\begin{equation}\label{eq:volume-variation}
 \left.\frac{d}{dt}\right|_{t=0}\Vol(X+tY)
 =\int_M\langle Y,\nu_X\rangle\,dA_X.
\end{equation}
For small $t$, the enclosed volume is a polynomial in $t$ of degree
at most three.
\end{lemma}
\begin{proof}
The assertion about embeddings follows from the openness just cited.
Orient $M$ so that $X$ has outward normal. On $\R^3$ define the
two-form
\[
 \beta=\tfrac13(x\,dy\wedge dz+y\,dz\wedge dx+z\,dx\wedge dy),
 \qquad d\beta=dx\wedge dy\wedge dz.
\]
Set $X_t=X+tY$ and $\mathcal V(t)=\int_M X_t^*\beta$.
The divergence theorem gives $\mathcal V(0)=\Vol(X)>0$.
For every embedded $X_t$ it also gives
$\mathcal V(t)=\Vol(X_t)$ or $-\Vol(X_t)$, according to orientation.
Continuity and positivity at zero select the positive sign for
all sufficiently small $|t|$.
In positively oriented local coordinates,
\[
 X_t^*\beta=\tfrac13
       \det(X+tY,X_u+tY_u,X_v+tY_v)\,du\wedge dv.
\]
The integrand is a polynomial of degree at most three, proving
the polynomial assertion globally.

Extend the field $Y$ off $X(M)$ by a tubular projection and use
Cartan's formula for the derivative of pullback at zero.
The exact term has integral zero by Stokes' theorem on the
closed source. The remaining term is
\[
 \mathcal V'(0)=\int_M X^*(\iota_Y d\beta)
       =\int_M\langle Y,\nu_X\rangle\,dA_X.
\]
Here $\iota_Y$ denotes contraction with the extended velocity field;
its restriction to the surface is independent of the extension.
This proves \eqref{eq:volume-variation}. The classical ingredients
are Cartan's formula, Stokes' theorem, and the divergence theorem;
see \cite[Theorem~14.35 and Chapter~16, including Theorem~16.32]{Lee}.
\end{proof}

\begin{proof}[Proof of Proposition~\ref{prop:pair}]
For arbitrary tangent vectors $v,w$, expand
\[
 \langle d(X+tY)(v),d(X+tY)(w)\rangle.
\]
The linear term vanishes by \eqref{eq:strain}, leaving
$g_X(v,w)+t^2\langle dY(v),dY(w)\rangle$. This proves the exact
identity \eqref{eq:exact-metric} at $t=\pm a$.
By Lemma~\ref{lem:volume},
\begin{equation}\label{eq:odd-volume}
 \Vol(X+aY)-\Vol(X-aY)
 =2a\int_M\langle Y,\nu_X\rangle\,dA_X+O(a^3).
\end{equation}
The coefficient of $a$ is nonzero; hence the difference is nonzero
for every sufficiently small $a>0$.

Choose a compact neighborhood of $\supp Y$ inside $\{K_X<0\}$.
For fixed $Y$, $X+tY\to X$ in $C^2$, so its curvature stays negative
on that neighborhood for small $|t|$. Outside it the embedding is
unchanged. Lemma~\ref{lem:tight-preservation} proves tightness.
Euclidean isometries preserve enclosed volume, so unequal volumes
exclude congruence of the images.
\end{proof}

\section{The BGQR annulus and completion to a tight torus}
\label{sec:seed}

The purpose of this section is only to prove Theorem~\ref{thm:seed}.
The local bending theorem will use the resulting torus and its one
retained curve, without any formula for that curve.
For the planar boundary curves below, \emph{strictly convex} means
that the curve is an embedded circle bounding a convex disk and
has positive signed curvature in its counterclockwise orientation.

\begin{definition}\label{def:planar-jets}
Let $S\subset\R^2\times[-h,h]$, $h>0$, be an embedded compact
annulus whose boundary curves bound smooth planar disks
$D_\pm\subset\{z=\pm h\}$. An \emph{inward projected collar} at the
upper boundary means that horizontal projection maps a neighborhood
of that boundary in $S$ diffeomorphically onto a collar of
$\partial D_+$ inside $D_+$. On this collar the annulus is a graph
$z=h+b_+(x)$. Infinite-order planar contact means that every partial
derivative of $b_+$, including its value, is zero on $\partial D_+$.
At the lower boundary the analogous graph is $z=-h+b_-(x)$ with
the same condition. Smoothness and derivatives at the boundary
are understood in coordinates on a manifold with boundary.
\end{definition}

\begin{proposition}[An annulus implies a tight completion]
\label{prop:completion}
Suppose an embedded smooth compact annulus $S$ has all the properties
in Definition~\ref{def:planar-jets}, its interior lies in $-h<z<h$,
and $K_S\le0$. Then it is contained in a smooth embedded tight torus
which agrees with $S$ throughout its interior.
\end{proposition}
\begin{proof}
Choose $R>0$ such that the two disks and the horizontal projection
of $S$ lie in $|x|<R$. To supply a convex outer surface with planar
caps, choose a smooth function $\phi:\R\to[0,\infty)$ satisfying
\[
 \phi=0\text{ on }(-\infty,0],\qquad
 \phi'>0,\ \phi''>0\text{ on }(0,\infty),\qquad
 \phi(t)\longrightarrow\infty.
\]
For example, for $t>0$ set
$\phi(t)=\int_0^t(t-u)e^{-1/u}\,du$ and extend by zero.
Consider the body
\begin{equation}\label{eq:body}
 \Omega=\{(x,z)\in\R^2\times\R:
                 \phi(|x|^2-R^2)+z^2/h^2\le1\}.
\end{equation}
Its defining function is convex, it has a minimum zero at the origin,
and the sublevel set is compact with nonempty interior.
On level one the gradient cannot vanish: if $z\ne0$ its vertical
component is nonzero, while if $z=0$, then $|x|>R$ and its horizontal
component is $2\phi' x\ne0$.
Thus $\partial\Omega$ is smooth. A compact convex body with nonempty
interior has boundary homeomorphic to a sphere; here radial projection
from the origin also gives a diffeomorphism, by transversality and
the implicit function theorem.

The two planar caps are $|x|\le R$, $z=\pm h$.
Away from them, $|x|>R$, and the defining Hessian is
\[
 \begin{pmatrix}
  2\phi'I+4\phi''x\otimes x&0\\0&2/h^2
 \end{pmatrix},
\]
which is positive definite. Its restriction to the tangent plane,
divided by the length of the gradient, is the second fundamental
form up to sign. Thus $\partial\Omega$ has positive Gaussian curvature
away from the caps and zero curvature on the caps.
Gauss--Bonnet on this convex sphere gives total curvature $4\pi$.

Every interior point of $S$ lies strictly inside $\Omega$: its
horizontal term in \eqref{eq:body} is zero and $|z|<h$.
Remove the interiors of $D_+$ and $D_-$ from $\partial\Omega$
and attach $S$ along their boundary curves. In horizontal coordinates,
the planar piece occupies the exterior side of each disk boundary,
and the annulus occupies the interior side. The graph defining the
union is $b_\pm$ on one side and zero on the other; the vanishing
of all boundary derivatives proves that this graph is smooth.
This also proves immersion along the seam. Strict interiority
excludes intersections between the other points of $S$ and
$\partial\Omega$.

The union is a connected closed embedded surface. It is orientable,
as every closed embedded surface in $\R^3$ is. Removing two disks
from a sphere leaves an annulus, so its Euler characteristic after
attaching $S$ is zero. The classification of closed orientable
surfaces therefore identifies it with a torus. The removed disks
have zero curvature, the retained convex part contributes positive
curvature $4\pi$, and $S$ contributes none. Its total positive
curvature is consequently $4\pi$.
Theorem~\ref{thm:tightness} proves tightness.
\end{proof}

\begin{theorem}[External construction: BGQR]\label{thm:bgqr}
There exist a smooth embedded compact annulus $S_0$, negatively curved
in its interior, and a regular embedded closed asymptotic curve
$\Gamma$ in that interior. The annulus has a parametrization
\[
 S_0(\theta,r)=(C(\theta)-r e_\theta,z(r)),
 \qquad (\theta,r)\in(\R/2\pi\mathbb Z)\times[-200,200],
\]
where $e_\theta=(\cos\theta,\sin\theta)$ and $J(x,y)=(-y,x)$.
There are periodic functions $P,Q$ such that
\[
 C'=P e_\theta+QJe_\theta,\quad P>0.
\]
The function $z$ is odd, $z(r)=r$ for $|r|\le175$, $z'>0$ for
$|r|<200$, $rz''\le0$, and
\[
 z(\pm200)=\pm575/3,\qquad z'(\pm200)=0,\qquad
 z''(\pm200)=\mp1/15.
\]
Writing $v=Q-r$ and $D=P^2+v^2+PQ'-vP'$, the end regions
$175\le|r|\le200$ satisfy $D>0$ and $rv<0$.
The endpoint sections are strictly convex planar Jordan curves.
\end{theorem}

This is an external input: Theorem~1.1, Section~2.4,
Proposition~3.2, and Proposition~3.3 of \cite{BGQR}, specifically
equation~(7) in Section~3.1 and equations~(11), (13), and~(14) in
Section~3.2. The theorem there supplies two curves; we select one.
The numerical and coordinate assertions above are outputs of that
particular construction, rather than consequences of its abstract
Theorem~1.1 alone. We do not reproduce its construction.
The source boundary contact is not infinite-order planar contact,
so Proposition~\ref{prop:completion} still requires the following
modification.

\begin{lemma}[Flattening the ends while retaining the curve]
\label{lem:flatten}
The annulus in Theorem~\ref{thm:bgqr} can be modified outside a
neighborhood of $\Gamma$ so that it satisfies all the hypotheses of
Proposition~\ref{prop:completion}, with $K<0$ in its interior.
\end{lemma}
\begin{proof}
Put $h_0=575/3$. Since $\Gamma$ is compactly contained in the
interior and $z$ is strictly increasing there, choose
\[
 \max\{175,\max_\Gamma|z|\}<h_1<h_0.
\]
Choose an even smooth function $q:\R\to[0,1]$ which equals one on
$[-h_1,h_1]$, is positive on $(-h_0,h_0)$, vanishes outside this
interval, and is nonincreasing on $[0,\infty)$.
To construct it, take the nonnegative bump $e^{-1/[t(1-t)]}$ on
$(0,1)$, extended by zero, and normalize its cumulative integral.
Rescaling this smooth step function and subtracting from one gives
the required transition on $[h_1,h_0]$; reflect evenly.
In particular all its derivatives vanish at $\pm h_0$.
Set
\[
 f(t)=\int_0^t q(a)\,da,\qquad H(r)=f(z(r)),\qquad h=f(h_0)>0.
\]
Then $f$ is odd, equals the identity on $[-h_1,h_1]$,
$f'>0$ in $(-h_0,h_0)$, and $tf''(t)\le0$.
The new annulus is
\[
 S(\theta,r)=(C(\theta)-re_\theta,H(r)).
\]
It agrees with $S_0$ in a neighborhood of $\Gamma$.
On the interior, $H'>0$; hence equality of two points of $S$
first gives equality of $r$, and then injectivity follows from
that of $S_0$. Its differential has rank two there. At an endpoint
the horizontal derivatives have determinant $v\ne0$, so immersion
also holds at the boundary. The height separates the two endpoint
sections from all other sections. Thus $S$ is an embedding of
the compact annulus.

For completeness the curvature calculation is given explicitly.
In the orthonormal horizontal basis $(e_\theta,Je_\theta)$,
write $w=C-re_\theta$. Then
\[
 w_\theta=P e_\theta+vJe_\theta,\quad w_r=-e_\theta,\quad
 w_{\theta\theta}=(P'-v)e_\theta+(P+Q')Je_\theta .
\]
The area vector of $S$ is
\[
 S_\theta\times S_r=(H'v,-H'P,v)
\]
in the basis $(e_\theta,Je_\theta,e_z)$, and its squared length is
$H'^2(P^2+v^2)+v^2$.
The unnormalized second coefficients are
\[
 \langle S_{\theta\theta},S_\theta\times S_r\rangle=-H'D,
 \quad
 \langle S_{\theta r},S_\theta\times S_r\rangle=H'P,
 \quad
 \langle S_{rr},S_\theta\times S_r\rangle=H''v .
\]
Consequently
\begin{equation}\label{eq:height-curvature}
 K_S=-\frac{H'^2P^2+H'H''vD}
             {[H'^2(P^2+v^2)+v^2]^2}.
\end{equation}
Where the height changes, $|z|>h_1>175$ implies $|r|>175$.
Since $z$ has the sign of $r$,
\[
 rH''=r\{f''(z)z'^2+f'(z)z''\}\le0.
\]
The inequalities $rv<0$ and $D>0$ therefore give $H''vD\ge0$.
The first numerator term in \eqref{eq:height-curvature} is strictly
positive in the interior because $H'>0$ and $P>0$.
Elsewhere the original surface is unchanged. Thus $K_S<0$
throughout the interior and $K_S=0$ on the boundary.

It remains to verify the attachment conditions, including the side
of the projected collars. The horizontal endpoint curves have
tangent $Pe_\theta+vJe_\theta$ and determinant $D>0$, so their
counterclockwise orientation has outward unit normal
\[
 \frac{ve_\theta-PJe_\theta}{\sqrt{P^2+v^2}}.
\]
At the upper end the continuation towards increasing $r$ is
$-e_\theta$ and its dot product with this outward normal is
$-v/\sqrt{P^2+v^2}>0$.
At the lower end the continuation towards decreasing $r$ is
$e_\theta$ and the corresponding dot product is
$v/\sqrt{P^2+v^2}>0$.
The interior projected collars therefore lie inside the bounded
disks. Their projection is a local diffeomorphism, since its
Jacobian is $v\ne0$ at the boundary. Compactness of each boundary
and its embedding provide a collar on which this projection is a
diffeomorphism.

All positive-order derivatives of $f$ at $\pm h_0$ vanish.
The chain rule shows that $H-H(\pm200)$ has every derivative
zero at the corresponding endpoint. After the smooth horizontal
change of coordinates just established, the graph functions also
have every boundary derivative zero. These are exactly the
conditions of Definition~\ref{def:planar-jets}.
\end{proof}

\begin{proof}[Proof of Theorem~\ref{thm:seed}]
Apply Lemma~\ref{lem:flatten} to the annulus and chosen curve in
Theorem~\ref{thm:bgqr}, and then apply
Proposition~\ref{prop:completion}. The curve and a neighborhood of
it were retained unchanged, so it is still embedded, asymptotic,
and in $K<0$ on the completed tight torus. Parametrize that torus
by $T^2$; the inverse parametrization of the retained curve gives
the required embedding $c$.
\end{proof}

\section{The Darboux frame and the forced inflection}\label{sec:frame}

The remaining sections prove Theorem~\ref{thm:one-curve} for any
starting torus satisfying its hypotheses. The construction of
Section~\ref{sec:seed} will not be used again.

Let $\gamma:\R/\ell\mathbb Z\to\R^3$, $\ell>0$, be the given curve
parametrized by arclength, and let $n(s)$ be the unit normal of
the surface along it. Put $T=\gamma_s$ and $E=n\times T$.
Then $(T,E,n)$ is a positively oriented orthonormal frame.
All derivatives in $s$ are derivatives of periodic lifts.

\begin{lemma}[Frame equations]\label{lem:frame}
There are smooth real periodic functions $k,\tau$ satisfying
\begin{equation}\label{eq:darboux}
 T_s=kE,\qquad E_s=-kT+\tau n,\qquad n_s=-\tau E.
\end{equation}
They are given by $k=\langle T_s,E\rangle$ and
$\tau=\langle E_s,n\rangle$. Moreover $K=-\tau^2$ along $\gamma$.
\end{lemma}
\begin{proof}
Unit speed gives $T_s\perp T$, and asymptoticity gives
$T_s\perp n$, proving the first equation.
Differentiating orthonormality gives the other two, with the
$T$ component of $n_s$ equal to
$-\langle n,T_s\rangle=0$.
In the orthonormal tangent basis $(T,E)$, the second fundamental
form has first entry zero and mixed entry $\tau$.
Its determinant is therefore $-\tau^2$, which proves the curvature
identity by \eqref{eq:gauss-curvature}.
\end{proof}

Here $k$ is the signed geodesic curvature and $\tau$ is the signed
geodesic torsion. Negative Gaussian curvature implies $\tau\ne0$;
hence its sign is constant on the circle. Define
\begin{equation}\label{eq:rho}
 \rho=|\tau|^{-1/2}>0.
\end{equation}
This smooth periodic function equals the negative Gaussian curvature
along the core raised to the power $-1/4$.
We use this surface frame even
at points where the curve's curvature vanishes, where a Frenet
normal need not be defined.
The exponent is chosen so that
$\rho_s/\rho=-\tau_s/(2\tau)$, the integrating factor relation
in the asymptotic return equation of Section~\ref{sec:ruled}.

\begin{definition}\label{def:inflection}
An \emph{inflection} of the unit-speed curve $\gamma$ is a point
$s_0$ where $\gamma_{ss}(s_0)=0$. By \eqref{eq:darboux} this is
equivalent to $k(s_0)=0$. No prescribed order of vanishing is
included in this definition.
\end{definition}

\begin{lemma}[Injectivity of the negative Gauss map]\label{lem:gauss}
For a smooth embedded tight torus with outward normal $n$, the
restriction of $n$ to $\{K<0\}$ is injective.
\end{lemma}
\begin{proof}
For a regular value $u\in S^2$ of $n$, let $m_+(u)$ and $m_-(u)$
be the numbers of preimages with positive and negative curvature.
They are finite because the source is compact and each preimage
is isolated. A maximum of the height function $\langle X,u\rangle$
has outward normal $u$. At that point its Hessian is negative
semidefinite; regularity of $u$ makes the Hessian nondegenerate,
so it is negative definite and the Gaussian curvature is positive.
Thus $m_+(u)\ge1$ for every regular value.

The change-of-variables formula for the Gauss map, whose absolute
Jacobian is $|K|$, and Theorem~\ref{thm:tightness} give
\[
 \int_{S^2}m_+(u)\,dA(u)=\int_{\{K>0\}}K\,dA=4\pi .
\]
Sard's theorem makes the set of nonregular values have measure zero.
Since $S^2$ has area $4\pi$, it follows that $m_+=1$ almost
everywhere. The signed-degree formula gives
\[
 m_+(u)-m_-(u)=\deg n
 =\frac1{4\pi}\int_{T^2}K\,dA=\frac{\chi(T^2)}2=0.
\]
Here the equality for the degree follows either from
$n^*(dA_{S^2})=K\,dA$ or the usual integration formula for degree,
and the next equality is Gauss--Bonnet. Thus $m_-=1$ almost
everywhere as well. Change of variables, Sard's theorem, and the
degree formula are used in their ordinary smooth-manifold forms;
see \cite[Chapters~6 and~17]{Lee}.

If two different points in $\{K<0\}$ had the same normal, the inverse
function theorem would provide disjoint neighborhoods mapped
diffeomorphically onto overlapping open subsets of $S^2$.
Every regular value in their nonempty open overlap would have at
least two negative-curvature preimages. This contradicts $m_-=1$
almost everywhere, since a nonempty open subset of $S^2$ has
positive area.
\end{proof}

\begin{proposition}[One loop supplies an inflection]\label{prop:inflection}
Every regular embedded closed asymptotic curve in $\{K<0\}$ on
a smooth embedded tight torus has an inflection. In fact its
signed geodesic curvature $k$ takes both positive and negative values.
\end{proposition}
\begin{proof}
Lemma~\ref{lem:gauss} shows that $s\mapsto n(s)$ is injective.
It is regular because $|n_s|=|\tau|>0$. It is therefore a simple
regular closed curve on the unit sphere.

We use the Gauss--Bonnet--Fenchel argument recorded by
Ghomi--Raffaelli \cite[Note~5.2, equations~(10)--(12)]{GR}, where
it is attributed to Kovaleva \cite{Kovaleva}. We give the calculation
to specify its orientations and its application here.
Let $\sigma$ be arclength on the spherical normal curve, and
let $\kappa_{\mathrm{sph}}$ be its signed geodesic curvature
for the outward orientation of $S^2$.
The frame equations give
\[
 n_{ss}=-\tau_sE+\tau kT-\tau^2n,\qquad
 n\times n_s=\tau T.
\]
Consequently
\[
 \kappa_{\mathrm{sph}}
 =\frac{\langle n_{ss},n\times n_s\rangle}{|n_s|^3}
 =\frac{k}{|\tau|},\qquad d\sigma=|\tau|\,ds,
\]
and therefore $\int\kappa_{\mathrm{sph}}\,d\sigma=\int_0^\ell k\,ds$.

Choose the disk on the sphere lying to the left of the oriented
normal curve, and let its area be $A$. The Jordan curve theorem
on the sphere gives two disks with nonempty open interiors, so
$0<A<4\pi$. The sphere has curvature one and this disk has Euler
characteristic one. Gauss--Bonnet with boundary gives
\[
 \int\kappa_{\mathrm{sph}}\,d\sigma+A=2\pi,
 \qquad
 \left|\int_0^\ell k\,ds\right|=|2\pi-A|<2\pi .
\]
Gauss--Bonnet with boundary is used in its standard smooth-boundary
form \cite[\S4-5]{doCarmo}. No convexity or hemisphere condition
on the spherical disk is required.

On the other hand Fenchel's theorem \cite{Fenchel} for a regular
closed space curve gives
\[
 \int_0^\ell|\gamma_{ss}|\,ds
 =\int_0^\ell|k|\,ds\ge2\pi.
\]
If $k$ had only one sign, including the possibility of zeros,
the last integral would equal $|\int k\,ds|$, a contradiction.
Thus $k$ has both signs, and continuity gives a zero.
\end{proof}

The role of the inflection is to permit a correction concentrated
near a point where $k=0$ without significantly changing the
transverse component of a ruling. Before choosing that ruling,
we arrange one further derivative condition. This is a modification
of the starting geometry, rather than an additional hypothesis on
the curve supplied by Theorem~\ref{thm:seed}.

\section{Local preparation of the curve}\label{sec:preparation}

\begin{proposition}[An inflection implies a prepared curve]
\label{prop:prepare}
Let $X_0:M\to\R^3$ be a smooth embedding of a closed oriented
surface, and let $c:\R/\ell\mathbb Z\to M$ be a smooth embedding
such that $\gamma=X_0\circ c$ is a unit-speed asymptotic curve
in $\{K_{X_0}<0\}$. Use its frame from
Section~\ref{sec:frame}, and suppose $k(s_0)=0$.
For every open neighborhood $U$ of $c(\R/\ell\mathbb Z)$ with
$U\subset\{K_{X_0}<0\}$ and every $\varepsilon>0$, there is a
smooth embedding $\widetilde X:M\to\R^3$ such that
\[
 \supp(\widetilde X-X_0)\Subset U,\quad
 \|\widetilde X-X_0\|_{C^2}<\varepsilon,\quad
 K_{\widetilde X}<0\text{ on }U.
\]
The actual curve $\widetilde X\circ c$ is closed, embedded, and
asymptotic. In its own arclength $\sigma$, chosen to be zero at
$c(s_0)$, its frame functions satisfy
\begin{equation}\label{eq:prepared}
 \widetilde k(0)=0,\qquad
 \widetilde\rho_{\sigma\sigma}(0)\ne\widetilde\rho(0)^{-3}.
\end{equation}
Here $\widetilde\rho=|\widetilde\tau|^{-1/2}$ for the new curve.
If $X_0$ is tight, $\widetilde X$ is tight.
\end{proposition}
\begin{proof}
We change the actual curve while retaining its tangent directions
and its normal directions. Prescribe a new velocity $v(s)T(s)$
with $v>0$. To keep the curve closed it is necessary and sufficient
that
\begin{equation}\label{eq:closure-moment}
 \int_0^\ell(v-1)T\,ds=0,
\end{equation}
because $\int_0^\ell T\,ds=0$. We construct a family $v=v_\delta$,
$\delta\downarrow0$, with
\begin{equation}\label{eq:speed-jets}
 v(s_0)=1,\quad v_s(s_0)=0,\quad v_{ss}(s_0)=d,
 \qquad \|v-1\|_{C^1}\longrightarrow0,
\end{equation}
where the fixed real number $d$ will be chosen below.

There is an interval disjoint from $s_0$ on which $k\ne0$.
Indeed $k\equiv0$ would make $T$ constant, contradicting
$\int T\,ds=0$ and $|T|=1$; a nonzero value has an open
neighborhood away from the zero at $s_0$. On such an interval,
\[
 T_s=kE,\quad T_{ss}=-k^2T+k_sE+k\tau n,\quad
 \det(T,T_s,T_{ss})=k^2\tau\ne0.
\]
It follows that the vectors $T(s)$ on that interval span $\R^3$:
otherwise they and their first two derivatives would lie in a
fixed proper linear subspace, contradicting the determinant.
Choose three points with linearly independent tangent vectors,
and choose smooth bumps $b_1,b_2,b_3$ of integral one in
disjoint sufficiently small neighborhoods of those points.
The three vectors $\int b_iT\,ds$ are then linearly independent.

Choose a smooth cutoff $\chi$ on $\R$, equal to one near zero
and supported in $(-1,1)$, and set in a coordinate interval at $s_0$
\[
 h_\delta(s)=\tfrac d2(s-s_0)^2
                        \chi((s-s_0)/\delta),
\]
extended by zero periodically. Its value and first derivative at
$s_0$ are zero, and its second derivative there is $d$.
Moreover $\|h_\delta\|_{C^1}=O(\delta)$ and
$|\int h_\delta T\,ds|=O(\delta^3)$.
Solve the invertible three-dimensional linear system
\[
 \sum_{i=1}^3 a_i\int b_iT\,ds=-\int h_\delta T\,ds.
\]
Its matrix is fixed, so $a_i=O(\delta^3)$.
Then $v=1+h_\delta+\sum_i a_i b_i$ satisfies
\eqref{eq:closure-moment} and \eqref{eq:speed-jets}.
For small $\delta$ it is positive.

Integrate $\widetilde\gamma_s=vT$ with
$\widetilde\gamma(s_0)=\gamma(s_0)$. The moment condition proves
periodicity; estimates on $v-1$ and $v_s$ show
$\widetilde\gamma-\gamma\to0$ in $C^2$.
Since the surface is oriented, its embedded source circle has a
two-sided smooth collar. Choose coordinates $(s,w)$ in such a
collar, with $X_0(s,0)=\gamma(s)$, whose closure lies in $U$.
For a fixed cutoff $\chi_0(w)$ equal to one near zero and
supported in the collar, define
\[
 \widetilde X(s,w)=X_0(s,w)+
              \chi_0(w)(\widetilde\gamma(s)-\gamma(s)),
\]
and put $\widetilde X=X_0$ elsewhere.
This tends to $X_0$ in $C^2$. Embedding and negative curvature
persist, and Lemma~\ref{lem:tight-preservation} preserves
tightness when present.

At $w=0$, the two tangent vectors of $\widetilde X$ are $vT$
and the original transverse tangent. Their span and orientation
are unchanged; hence the normal is exactly $n(s)$.
The new acceleration
$\widetilde\gamma_{ss}=v_sT+vkE$ is tangent to this plane,
so the new curve is asymptotic.
Let $\sigma$ be its actual arclength, with $\sigma(s_0)=0$.
Its new period is
$\widetilde\ell=\int_0^\ell v(s)\,ds$; vector closure does not
require this number to equal $\ell$.
Then
\[
 d\sigma=v\,ds,\qquad
 \widetilde k=k/v,\qquad
 \widetilde\tau=\tau/v,\qquad
 \widetilde\rho=\sqrt v\,\rho.
\]
At $s_0$, where $v=1$ and $v_s=0$, differentiating twice yields
\[
 \widetilde k=0,\qquad
 \widetilde\rho_{\sigma\sigma}=\rho_{ss}+\tfrac12\rho d,
 \qquad
 \widetilde\rho^{-3}-\widetilde\rho_{\sigma\sigma}
 =\rho^{-3}-\rho_{ss}-\tfrac12\rho d .
\]
Choose $d$ different from
$2(\rho^{-3}-\rho_{ss})(s_0)/\rho(s_0)$.
The last expression is then nonzero, proving
\eqref{eq:prepared} in the new arclength. A change of parameter
alone would not change this geometric condition; the construction
has changed the curve and the containing surface.
\end{proof}

\section{From a prepared curve to a periodic ruled annulus}
\label{sec:ruled}

The next statement only uses a framed space curve. It is independent
of any completion to a closed surface. Its output will be inserted
into the prepared surface in Section~\ref{sec:insertion}.

\begin{proposition}[Prepared curve implies periodic ruled annuli]
\label{prop:ruled}
Suppose $\gamma:\R/\ell\mathbb Z\to\R^3$ is a smooth embedded
unit-speed curve with a smooth positively oriented orthonormal frame
$(T,E,n)$ satisfying $T=\gamma_s$, \eqref{eq:darboux}, and
$\tau\ne0$. Put $\rho=|\tau|^{-1/2}$.
If for some $s_0$
\[
 k(s_0)=0,\qquad
 \rho(s_0)^{-3}-\rho_{ss}(s_0)\ne0,
\]
then there exist smooth periodic functions $f>0$, $g$, and $\omega$
such that, for $V=gT+fE$:
\begin{enumerate}
 \item $R(s,u)=\gamma(s)+uV(s)$ is an embedded negatively curved
 ruled strip on $|u|<a$ for some $a>0$, with the prescribed tangent
 plane $\operatorname{span}(T,E)$ along $u=0$;
 \item $\det(\gamma_s,V,V_s)=f^2\tau\ne0$;
 \item for all sufficiently large $c$,
 \[
 B(s,c)=\gamma(s)+\frac{\rho(s)}{c+\omega(s)}V(s)
 \]
 has closed asymptotic constant-$c$ curves.
 For every bounded nonempty open interval of such $c$, this is an embedding of
 the corresponding open annulus, and its image approaches the
 core as the interval is moved towards $+\infty$.
\end{enumerate}
\end{proposition}

\begin{proof}
Write a transverse vector in the prescribed tangent plane as
$V=gT+fE$, where $f>0$. Its scale is part of the choice: fixing
that scale in advance would impose an unnecessary nonlinear condition
on the return equation. We impose the linear relation
\begin{equation}\label{eq:ruling-constraint}
 f_s=-kg.
\end{equation}
For the moment $f,g$ are unspecified. Define the geometric component
$\xi=\langle V_s,T\rangle=g_s-kf$. Then
\[
 V_s=\xi T+f\tau n,\qquad
 V_{ss}=\xi_sT+(k\xi-f\tau^2)E+(f_s\tau+f\tau_s)n.
\]
In particular
\begin{equation}\label{eq:delta}
 \Delta:=\det(T,V,V_s)=f^2\tau\ne0.
\end{equation}
At $u=0$ the ruling is transverse to the core, so $R$ is an
immersion near it. Its mixed unnormalized second coefficient is
$\langle R_{su},R_s\times R_u\rangle=\Delta$ and $R_{uu}=0$.
Formula \eqref{eq:gauss-curvature} therefore gives
\begin{equation}\label{eq:ruled-curvature}
 K_R=-\frac{\Delta^2}{|R_s\times R_u|^4}<0
\end{equation}
where $R$ is immersed.

We determine the conditions for the nonruling asymptotic curves to
close. Their graphs are $s\mapsto R(s,u(s))$.
The acceleration has normal component equal to that of
$R_{ss}+2u_sR_{su}$: the $u_{ss}R_u$ term is tangent and
$R_{uu}=0$. Set
\[
 2A=\xi_s-\frac{\tau_s}{\tau}\xi+g\tau^2 .
\]
The frame equations give the two unnormalized coefficients
\[
 \langle R_{ss},R_s\times V\rangle
     =f^2\tau_su-2f^2\tau Au^2,\qquad
 \langle R_{su},R_s\times V\rangle=f^2\tau.
\]
Thus the asymptotic equation is the Riccati equation
\begin{equation}\label{eq:riccati}
 u_s=\frac{\rho_s}{\rho}u+Au^2.
\end{equation}
Here $\tau_s/\tau=-2\rho_s/\rho$ follows from
$\rho=|\tau|^{-1/2}$. For $u\ne0$, taking the reciprocal gives
\begin{equation}\label{eq:reciprocal}
 (\rho/u)_s=-\rho A.
\end{equation}
It follows that periodicity is equivalent to the single condition
$\int_0^\ell\rho A\,ds=0$. Indeed integration over a period proves
necessity; if this integral vanishes, a periodic function $\omega$
with $\omega_s=-\rho A$ exists, and every nonzero solution with
positive denominator is
\begin{equation}\label{eq:leaves}
 u(s,c)=\frac{\rho(s)}{c+\omega(s)}.
\end{equation}
The parameter $c$ is the additive constant in $\rho/u$; adding a
constant to $\omega$ only changes this label.

It remains to choose $f,g$ so that both
\eqref{eq:ruling-constraint} and the period condition hold.
The point of retaining the scale $f$ is that these requirements
become linear. Using $\tau^2=\rho^{-4}$ and
$g_s=\xi+kf$, direct differentiation gives
\begin{equation}\label{eq:period-identity}
 2\rho A-\{g(\rho^{-3}-\rho_{ss})-fk\rho_s\}
       =(\rho\xi+\rho_sg)_s.
\end{equation}
The derivative on the right has integral zero. Consequently it
suffices to solve
\begin{equation}\label{eq:linear-periods}
 f_s+kg=0,\qquad
 \int_0^\ell\{g(\rho^{-3}-\rho_{ss})-fk\rho_s\}\,ds=0,
 \qquad f>0.
\end{equation}

We now use the prepared inflection. Since a closed unit-speed curve
cannot have constant tangent, $\int_0^\ell k^2\,ds>0$.
Let $\varphi_\delta$ be a nonnegative smooth function of integral
one supported in the $\delta$-neighborhood of $s_0$ on the circle.
Define
\[
 b_\delta=\varphi_\delta-
       \frac{\int_0^\ell k\varphi_\delta\,ds}
            {\int_0^\ell k^2\,ds}\,k,\qquad
 (P_\delta)_s=kb_\delta,\qquad P_\delta(s_*)=0,
\]
where $s_*$ is any fixed base point.
The integral of $kb_\delta$ is zero, so $P_\delta$ is periodic.
Since $k(s_0)=0$,
\[
 \int_0^\ell|kb_\delta|\,ds\longrightarrow0,\quad
 \|P_\delta\|_\infty\longrightarrow0,\quad
 \int_0^\ell b_\delta F\,ds\longrightarrow F(s_0)
\]
for every fixed continuous periodic $F$.
The first assertion follows by bounding the bump contribution by
$\sup_{|s-s_0|<\delta}|k(s)|$ and the correction by a constant
times $|\int k\varphi_\delta|$; both tend to zero.
The second follows by integration, and the third follows from
concentration of $\varphi_\delta$ and the vanishing correction.

Set
\begin{equation}\label{eq:balance}
 \lambda_\delta=
 -\frac{\int_0^\ell k\rho_s\,ds}
 {\int_0^\ell\{b_\delta(\rho^{-3}-\rho_{ss})
                         +P_\delta k\rho_s\}\,ds},
 \qquad
 f=1+\lambda_\delta P_\delta,\quad g=-\lambda_\delta b_\delta .
\end{equation}
The denominator converges to
$\rho(s_0)^{-3}-\rho_{ss}(s_0)\ne0$.
Thus $\lambda_\delta$ is bounded for small $\delta$, and
$f\to1$ uniformly, giving $f>0$. The first equality in
\eqref{eq:linear-periods} follows by differentiating $f$.
The second follows by substitution of \eqref{eq:balance}.
Fix one such $\delta$ once and for all. This fixes the smooth
ruling $V$, all its derivatives, and a periodic $\omega$.

For completeness, compactness first gives a uniform width on
which $R$ is immersed. It is a local embedding at every $(s,0)$
by transversality. If no narrower strip were injective, there would
be pairs of distinct parameters with $|u|,|u'|\to0$ and equal
images. Subsequences of their circle coordinates converge.
Equality of the limiting core points and embeddedness of $\gamma$
force both limits to be the same $s$. A local embedding near
$(s,0)$ then excludes those pairs, a contradiction.
After narrowing once more, injectivity on the compact closed
cylinder gives a homeomorphism onto its image. Its restriction
to the interior is therefore an embedded strip of some width $a>0$.

Finally, for $c$ sufficiently large, $c+\omega>0$ and
$0<u(s,c)<a$, uniformly in $s$. The map
$(s,c)\mapsto(s,u(s,c))$ is a diffeomorphism onto its image because
\[
 u_c=-\frac{\rho}{(c+\omega)^2}<0
\]
and its inverse is $c=\rho(s)/u-\omega(s)$.
Composing with $R$ proves the asserted annular embedding.
Equation \eqref{eq:reciprocal} proves that its constant-$c$ circles
are asymptotic. Their distance from the core tends uniformly to
zero as $c\to+\infty$.
The width was chosen after the ruling; no bound uniform in
$\delta$ on derivatives of $V$ is required.
\end{proof}

\section{A reciprocal annulus implies a supported bending}
\label{sec:bending}

This section gives a separate local result whose input is the
annulus itself. In particular the result does not require the core
to be parametrized by arclength or the function $\rho$ to have
any prescribed curvature interpretation. In its application they
will be exactly the objects supplied by Proposition~\ref{prop:ruled}.

\begin{theorem}\label{thm:band-bending}
Let $\ell>0$, and let $\gamma,V:\R/\ell\mathbb Z\to\R^3$ and
$\rho,\omega:\R/\ell\mathbb Z\to\R$ be smooth periodic maps with
$\rho>0$ and
\[
 \Delta(s)=\det(\gamma_s,V,V_s)\ne0.
\]
Let $I$ be a nonempty open interval contained in
$(-\min\omega,\infty)$. Suppose
\begin{equation}\label{eq:abstract-band}
 B(s,c)=\gamma(s)+u(s,c)V(s),\qquad
 u(s,c)=\frac{\rho(s)}{c+\omega(s)}
\end{equation}
is an embedding on $(\R/\ell\mathbb Z)\times I$, and each
constant-$c$ circle is asymptotic on this annulus.
Then $B$ has a smooth infinitesimal bending $Y$ whose support
is compactly contained in this parameter annulus and whose signed
normal flux satisfies
\begin{equation}\label{eq:band-flux}
 \int_{(\R/\ell\mathbb Z)\times I}
       \det(Y,B_s,B_c)\,ds\,dc\ne0.
\end{equation}
\end{theorem}

The two coordinates are asymptotic: the constant-$c$ curves are
asymptotic by hypothesis, and the constant-$s$ curves lie on straight
rulings, even though $c$ is not an affine parameter on them.
We first solve the strain equations along these closed curves and
then choose the remaining one-variable function to avoid cancellation
of the flux.

\subsection{Solving all three strain equations}
\begin{proof}[Construction of the bending]
Let $F\in C_c^\infty(I)$, to be chosen later. The notation
$C_c^\infty(I)$ means the smooth real functions on $I$ whose support
is a compact subset of $I$. Define
\begin{equation}\label{eq:pairings}
 Q=\sqrt{|\Delta|}\,u_c,\qquad p=QF(c),\qquad
 Y\cdot V=0,\quad Y\cdot B_s=p,\quad
 Y\cdot V_s=\frac{p_c}{2u_c}.
\end{equation}
Here $u_c<0$, so $Q$ never vanishes. Also
$B_s=\gamma_s+u_sV+uV_s$, whence
$\det(B_s,V,V_s)=\Delta$.
The three pairing equations therefore determine exactly one smooth
field $Y$. The data are periodic in $s$, so $Y$ is periodic as
well. Its support is contained in
$(\R/\ell\mathbb Z)\times\supp F$ and hence is compact in the
open annulus. Derivatives of $F$ do not enlarge its support.

We verify \eqref{eq:strain} on the coordinate pairs $cc,sc,ss$.
Since $B_c=u_cV$ and $V$ is independent of $c$,
differentiation of $Y\cdot V=0$ gives
$Y_c\cdot B_c=0$. For the mixed equation, put
$z=p_c/(2u_c)$ temporarily. Differentiating the pairings gives
\[
 Y_s\cdot V=-z,\qquad
 Y_c\cdot B_s=p_c-Y\cdot B_{sc}=p_c-u_cz,
\]
because $B_{sc}=u_{cs}V+u_cV_s$.
Thus
$Y_s\cdot B_c+Y_c\cdot B_s=-u_cz+p_c-u_cz=0$.

For the $ss$ equation, asymptoticity means that $B_{ss}$ is tangent.
As $B_s,V$ are independent, write
$B_{ss}=\alpha B_s+\beta V$ with smooth coefficients.
Differentiating
$\det(B_{ss},B_s,V)=0$ in $c$ yields
\[
 2u_{cs}\Delta+
 u_c\det(B_s,V,V_{ss})-\alpha u_c\Delta=0.
\]
Differentiating $\Delta=\det(B_s,V,V_s)$ in $s$ yields
\[
 \Delta_s=\alpha\Delta+\det(B_s,V,V_{ss}).
\]
Eliminating the last determinant proves
\begin{equation}\label{eq:transport}
 \alpha=\frac{\Delta_s}{2\Delta}+\frac{u_{cs}}{u_c}
        =\frac{Q_s}{Q}.
\end{equation}
Since $F$ is constant in $s$, $p_s=\alpha p$.
It follows that
$Y_s\cdot B_s=p_s-Y\cdot B_{ss}=p_s-\alpha p=0$.
All three components of the symmetric strain vanish.
\end{proof}

\begin{remark}[The Codazzi interpretation of the factor]
\label{rem:codazzi}
The factor $Q$ in \eqref{eq:pairings} has a geometric meaning.
Let $\mathcal W=|B_s\times B_c|$ be the area density and
let $\sigma=\mathrm{II}_{sc}$ for the parameter unit normal.
Then
\[
 |\sigma|\mathcal W
   =|\det(B_{sc},B_s,B_c)|=|\Delta|u_c^2=Q^2.
\]
In asymptotic coordinates, the Mainardi--Codazzi equations give
$\partial_s\log|\sigma|=\Gamma^s_{ss}-\Gamma^c_{sc}$,
while the area-density identity gives
$\partial_s\log\mathcal W=\Gamma^s_{ss}+\Gamma^c_{sc}$.
Here the $\Gamma$'s are the Christoffel coefficients of the induced
metric in $(s,c)$ coordinates. Adding the two equalities shows
$\Gamma^s_{ss}=\frac12\partial_s\log(|\sigma|\mathcal W)$.
Since $B_{ss}=\alpha B_s+\beta V$, its coefficient
$\Gamma^s_{ss}$ is $\alpha$. This recovers
\eqref{eq:transport} from the classical theory
\cite[\S4-3]{doCarmo}. The preceding determinant computation
proves the precise identity needed without requiring any
separate bending theorem.
\end{remark}

\subsection{Computing and selecting the flux}
\begin{proof}[Completion of the proof of Theorem~\ref{thm:band-bending}]
We must show that some choice of $F$ in \eqref{eq:pairings} has
nonzero flux. Introduce the area vector for the affine ruling
parameter and its squared length:
\[
 N(s,u)=(\gamma_s+uV_s)\times V,\qquad S(s,u)=|N(s,u)|^2 .
\]
Then $B_s\times B_c=u_cN$.
Solving the pairing equations gives
\[
 Y=\frac{p(V\times V_s)+(p_c/(2u_c))N}{\Delta}.
\]
Indeed the three scalar products are those of \eqref{eq:pairings};
uniqueness therefore verifies the formula. Since
$N_u=V_s\times V$, one has
$(V\times V_s)\cdot N=-S_u/2$.
Consequently
\[
 \det(Y,B_s,B_c)
       =\frac{Sp_c-u_cS_up}{2\Delta}.
\]
As $\Delta$ is independent of $c$ and $S_c=u_cS_u$, subtracting
$\partial_c(Sp/(2\Delta))$ leaves $-u_cS_up/\Delta$.
Compact support of $F$ removes this total derivative under
integration. Thus
\begin{equation}\label{eq:flux-kernel}
 \int\det(Y,B_s,B_c)\,ds\,dc
 =-\operatorname{sign}(\Delta)\int_I F(c)H(c)\,dc,\qquad
 H(c)=\int_0^\ell\frac{u_c^2S_u}{\sqrt{|\Delta|}}\,ds.
\end{equation}
The sign of $\Delta$ is constant since it is nonzero on a connected
circle.

The function $N$ is affine in $u$, so write
$S=S_0+S_1u+S_2u^2$, where the coefficients are smooth periodic
functions of $s$ and
$S_2=|V_s\times V|^2>0$ by $\Delta\ne0$.
Substituting $u$ and $u_c$ gives
\begin{equation}\label{eq:H}
 H(c)=\int_0^\ell\frac1{\sqrt{|\Delta|}}
 \left\{\frac{\rho^2S_1}{(c+\omega)^4}
       +\frac{2\rho^3S_2}{(c+\omega)^5}\right\}\,ds.
\end{equation}
Although the annulus need only be embedded for $c\in I$,
this scalar expression is defined on the whole connected interval
$c>-\min\omega$.

It is real analytic on that interval. To see this, fix $c_0$ there
and let $a=\min_s(c_0+\omega(s))>0$.
For $|c-c_0|<a/2$, the binomial power series for both denominator
powers in \eqref{eq:H} converge absolutely and uniformly in $s$,
with a summable bound on every smaller parameter interval.
Termwise integration gives a convergent power series for $H$.
Only analyticity in the parameter $c$ is used; the geometric
data need only be smooth.

We next prove that $H$ is not identically zero.
Set $\epsilon=c+\min\omega>0$ and
$r(s)=\omega(s)-\min\omega\ge0$. The integrand in \eqref{eq:H}
can be written
\[
 \frac{a(s)(\epsilon+r(s))+b(s)}{(\epsilon+r(s))^5},
 \qquad
 a(s)=\frac{\rho^2S_1}{\sqrt{|\Delta|}},\quad
 b(s)=\frac{2\rho^3S_2}{\sqrt{|\Delta|}}.
\]
These functions are bounded, and $m=\min_s b(s)>0$.
Let $M=\max_s|a(s)|$. Choose $r_0>0$ with
$2Mr_0\le m/2$, with any $r_0>0$ if $M=0$.
For $0<\epsilon\le r_0$ and $r(s)\le r_0$, the numerator is at
least $m/2$. On $\{r\ge r_0\}$ the integrand is uniformly bounded
in absolute value as $\epsilon\downarrow0$.

Choose any minimum point $s_*$. Smoothness gives $r_s(s_*)=0$
and $r(s_*+t)\le Ct^2$ in a coordinate interval, for a fixed
$C>0$. For small $\epsilon$, the interval $|t|\le\sqrt\epsilon$
lies in $\{r\le r_0\}$ and
$\epsilon+r(s_*+t)\le(1+C)\epsilon$.
Its integral is therefore bounded below by a positive constant
times $\epsilon^{-9/2}$. All the rest of $\{r\le r_0\}$
contributes nonnegatively, and the remaining region contributes
a uniformly bounded quantity. We have proved
\[
 H(c)\longrightarrow+\infty\quad
                 \text{as }c\downarrow-\min\omega.
\]
This argument applies to degenerate, nonisolated, or multiple minima.

By the identity theorem for real analytic functions, $H$ cannot
vanish identically on any nonempty open subinterval of its domain.
Choose a nonzero nonnegative $\chi\in C_c^\infty(I)$ and set
$F=\chi H$. Then $F\in C_c^\infty(I)$ and
\[
 \int_I F(c)H(c)\,dc=\int_I\chi(c)H(c)^2\,dc>0.
\]
Indeed $\chi>0$ on some open interval, on which $H$ cannot
vanish identically. Equation \eqref{eq:flux-kernel} now proves
\eqref{eq:band-flux}. The extension of the scalar formula outside
$I$ was used only to prove this nonidentity. The bending itself
is defined and supported entirely on the actual embedded annulus.
\end{proof}

\section{Insertion and passage to a global bending}
\label{sec:insertion}

The ruled annulus just constructed has the correct tangent planes
along its core, but need not already be part of the original surface.
The next proposition places that same ruled strip in the surface.
Its curvature estimate is needed because cutting off a quadratic
normal displacement need not be small in $C^2$.

\begin{proposition}[A tangent ruled strip can be inserted]
\label{prop:insert}
Let $X_0:M\to\R^3$ be a smooth embedding of a closed oriented
surface. Let $c:\R/\ell\mathbb Z\to M$ be an embedded asymptotic
curve with image $\gamma=X_0\circ c$, parametrized by arclength,
and let $U$ be an open source neighborhood contained in
$\{K_{X_0}<0\}$.
Suppose $R(s,u)=\gamma(s)+uV(s)$ is an embedded strip near $u=0$,
with $V(s)$ tangent to $X_0(M)$ and transverse to $\gamma_s$.
For every $\varepsilon>0$ there is a smooth embedding
$\widetilde X:M\to\R^3$ such that
\[
 \supp(\widetilde X-X_0)\Subset U,\qquad
 \|\widetilde X-X_0\|_{C^1}<\varepsilon,\qquad
 K_{\widetilde X}<0\text{ on }U,
\]
and there exist $a>0$ and a smooth embedding
$q:(\R/\ell\mathbb Z)\times(-a,a)\to U$ with
\[
 q(s,0)=c(s),\qquad \widetilde X(q(s,u))=R(s,u).
\]
If $X_0$ is tight, $\widetilde X$ is tight.
\end{proposition}
\begin{proof}
Use an ambient tubular neighborhood of $X_0(M)$ with its normal
projection $\pi$ onto $X_0(M)$; see
\cite[Theorem~6.24]{Lee}. After narrowing the strip, define
\[
 q=X_0^{-1}\circ\pi\circ R .
\]
At $u=0$, $q(s,0)=c(s)$, and $dq$ is invertible because
$R_s,R_u$ span the original tangent plane.
Thus $q$ is a local diffeomorphism near the core. It is a
diffeomorphism onto a source neighborhood for a sufficiently
narrow uniform strip. The proof of uniform injectivity is the same
compactness argument as in Proposition~\ref{prop:ruled}, now
applied to $q$ and its embedded core.
By the definition of normal projection, in that neighborhood
\[
 R\circ q^{-1}=X_0+hn
\]
for a smooth scalar function $h$ and the original surface normal
$n$. Equality of points and tangent planes along the core gives
$h=dh=0$ there.

Fix an original-source collar $(s,w)$ of the core with compact
closure in $U$. Narrow it to lie in the image of $q$. Here $s$
agrees with the arclength coordinate on $w=0$; it is not being
identified with the first parameter of $q^{-1}$ away from the core.
Taylor's formula with integral remainder gives
\[
 h(s,w)=w^2b(s,w)
\]
with $b$ smooth. Take a fixed smooth cutoff $\chi$ equal to one
on $[-1,1]$ and zero outside $(-2,2)$. For small $\eta>0$ put
\[
 h_\eta(s,w)=\chi(w/\eta)w^2b(s,w),\qquad
 X_\eta=X_0+h_\eta n
\]
on the collar and $X_\eta=X_0$ elsewhere.
For $\eta$ small the support lies compactly in $U$.
All the following estimates are uniform on the collar, with constants
depending on the already fixed strip:
\[
 h_\eta=O(\eta^2),\quad
 (h_\eta)_s=O(\eta^2),\quad(h_\eta)_w=O(\eta),
\]
\[
 (h_\eta)_{ss}=O(\eta^2),\quad
 (h_\eta)_{sw}=O(\eta),\quad(h_\eta)_{ww}=O(1).
\]
Thus $X_\eta\to X_0$ in $C^1$, giving immersion and global
embedding for small $\eta$. Its consistently chosen unit normal
satisfies $n_\eta-n=O(\eta)$.

Let $L_\eta,M_\eta,N_\eta$ be the second fundamental coefficients
of $X_\eta$ in the fixed collar. For the original surface
$L_0(s,0)=0$ by asymptoticity, so $L_0(s,w)=O(w)$.
Also $M_0(s,0)$ is nowhere zero, because at the core the
determinant of the second fundamental form is
$-M_0(s,0)^2<0$. The derivative estimates just stated and
$n_\eta-n=O(\eta)$ imply on $|w|\le2\eta$
\[
 L_\eta=O(\eta),\qquad
 M_\eta=M_0(s,0)+O(\eta),\qquad N_\eta=O(1).
\]
For example the $ss$ derivative of the added normal displacement
is $O(\eta^2)$, and the change of normal in the bounded original
$ss$ coefficient contributes $O(\eta)$; the mixed and $ww$
estimates follow in the same way. Hence
\[
 L_\eta N_\eta-M_\eta^2
     =-M_0(s,0)^2+O(\eta)<0
\]
uniformly for small $\eta$, since $\min_s|M_0(s,0)|>0$.
The first fundamental determinant stays bounded below by a
positive constant by $C^1$ convergence. Formula
\eqref{eq:gauss-curvature} proves $K_{X_\eta}<0$ on the
modified collar. Outside it the surface is unchanged, so the
curvature is negative on all of $U$.

Fix $\eta$ small enough for all the preceding conclusions and for
$\|X_\eta-X_0\|_{C^1}<\varepsilon$.
On $|w|<\eta$, $X_\eta=X_0+hn=R\circ q^{-1}$ exactly.
By compactness of the core, there is $a>0$ such that the image
under $q$ of $|u|<a$ lies in $|w|<\eta$. Restricting $q$ to
this strip gives exactly the inclusion asserted in the statement.
If $X_0$ is tight, Lemma~\ref{lem:tight-preservation} applies.
\end{proof}

\begin{lemma}[An actual supported band gives a global bending]
\label{lem:globalize}
Let $X:M\to\R^3$ be an embedding of a closed surface, let
$\ell>0$, let $I\subset\R$ be a nonempty open interval, and put
$A=(\R/\ell\mathbb Z)\times I$. Let $\iota:A\to M$ be a smooth
annular embedding with $X\circ\iota=B$.
If $B$ has an infinitesimal bending $Y_A$ compactly supported in $A$
and with $\int_A\det(Y_A,B_s,B_c)\,ds\,dc\ne0$, then $X$ has an
infinitesimal bending $Y$ supported in $\iota(\supp Y_A)$ and
with nonzero outward normal flux.
\end{lemma}
\begin{proof}
An embedding of equal-dimensional manifolds has open image, by
the inverse function theorem. Define $Y\circ\iota=Y_A$ on that
image and $Y=0$ outside it.
The image of $\supp Y_A$ is a compact subset of the open set
$\iota(A)$. Every point outside $\iota(A)$ therefore has a
neighborhood disjoint from this image, on which the extended
field is zero. Thus the extension is smooth.
The chain rule and the literal identity $X\circ\iota=B$ transfer
\eqref{eq:strain} to $\iota(A)$; outside the support it is zero.
On this connected annulus the outward normal of $X$ is either
the parameter normal $(B_s\times B_c)/|B_s\times B_c|$
or its negative, with a single constant sign. The
change-of-variables formula therefore gives
\[
 \int_M\langle Y,\nu_X\rangle\,dA_X
      =\pm\int_A\det(Y_A,B_s,B_c)\,ds\,dc\ne0.
\]
\end{proof}

\begin{proof}[Proof of Theorem~\ref{thm:one-curve}]
Choose a collar with compact closure in $U$.
Proposition~\ref{prop:inflection} supplies an inflection on the
original curve. Apply Proposition~\ref{prop:prepare}, with
$C^2$ error as small as needed, supported in that collar.
Call the resulting tight embedding $X_1$.
Parametrize its actual new core by its own arclength and equip
it with its Darboux frame. It satisfies precisely the hypotheses
of Proposition~\ref{prop:ruled}.

Fix the ruling supplied by that proposition.
Apply Proposition~\ref{prop:insert} to $X_1$ and this ruled strip,
again with support compactly in $U$, to obtain $X$.
Choose the errors in the two modifications so that their sum
in $C^1$ is less than the prescribed $\varepsilon$.
The surface remains negatively curved on all of $U$ and is tight.
The inclusion of a neighborhood of the ruled core is exact.

Let $q$ and $a>0$ be supplied by Proposition~\ref{prop:insert}.
Only now choose a bounded nonempty open interval $I$ of sufficiently
large leaf parameters $c$ so that $c+\omega(s)>0$ and
$0<u(s,c)<a$ for every $(s,c)\in(\R/\ell\mathbb Z)\times I$.
The annulus
$B(s,c)=\gamma(s)+\rho(s)V(s)/(c+\omega(s))$
lies inside that exact neighborhood.
Set $\iota(s,c)=q(s,u(s,c))$.
Then $\iota$ is an annular embedding and
$X\circ\iota=B$. Theorem~\ref{thm:band-bending} supplies a field
on this same annulus with compact support and nonzero signed
flux. Lemma~\ref{lem:globalize} produces the desired global
field $Y$. Its compact support lies in $U$.

This proves every assertion in \eqref{eq:one-curve-output}.
The choices were made in the order dictated by their dependencies:
the local preparation, then the ruling, then the insertion width,
then the leaf interval, and finally the bending function.
No limiting assertion about an earlier choice is used after a
later choice has been fixed.
\end{proof}

\section*{Acknowledgements}
I thank my mentor Xiang Ma for helpful discussions and encouragement.
AI-assisted tools were used in the preparation of this manuscript.

\clearpage
\begin{thebibliography}{99}
\bibitem{Aleksandrov}
V.~A.~Aleksandrov,
\emph{Sabitov's conjecture that volume is stationary under infinitesimal
bending of a surface}, Siberian Math. J. \textbf{30} (1989), 678--684,
\href{https://doi.org/10.1007/BF00971257}{doi:10.1007/BF00971257}.
\bibitem{BK}
T.~F.~Banchoff and W.~K\"uhnel,
\emph{Tight submanifolds, smooth and polyhedral}, in
\emph{Tight and Taut Submanifolds}, MSRI Publications \textbf{32},
Cambridge University Press, 1997, 51--118,
\href{https://library.slmath.org/books/Book32/files/banchoff.pdf}{available from SLMath}.
\bibitem{BGQR}
H.~Banks, M.~Ghomi, H.~Qiu, and M.~Raffaelli,
\emph{A counterexample to Nirenberg's conjecture on closed asymptotic
curves}, \href{https://arxiv.org/abs/2610.13848v2}{arXiv:2610.13848v2},
21 October 2026.
\bibitem{doCarmo}
M.~P.~do Carmo, \emph{Differential Geometry of Curves and Surfaces},
Prentice-Hall, 1976.
\bibitem{Fenchel}
W.~Fenchel, \emph{\"Uber Kr\"ummung und Windung geschlossener Raumkurven},
Math. Ann. \textbf{101} (1929), 238--252,
\href{https://doi.org/10.1007/BF01454836}{doi:10.1007/BF01454836}.
\bibitem{GR}
M.~Ghomi and M.~Raffaelli,
\emph{Topology of closed asymptotic curves on negatively curved surfaces},
J. Geom. Anal. \textbf{35} (2025), article 381,
\href{https://doi.org/10.1007/s12220-025-02200-3}{doi:10.1007/s12220-025-02200-3};
\href{https://arxiv.org/abs/2412.19266v2}{arXiv:2412.19266v2}.
\bibitem{Hirsch}
M.~W.~Hirsch, \emph{Differential Topology},
Graduate Texts in Mathematics \textbf{33}, Springer, 1976,
\href{https://doi.org/10.1007/978-1-4684-9449-5}{doi:10.1007/978-1-4684-9449-5}.
\bibitem{Kovaleva}
G.~A.~Kovaleva,
\emph{Absence of closed asymptotic lines on tubes of negative Gaussian
curvature with a one-to-one spherical mapping},
Fundam. Prikl. Mat. \textbf{1} (1995), no.~4, 953--977 (Russian).
\bibitem{Lee}
J.~M.~Lee, \emph{Introduction to Smooth Manifolds}, second edition,
Graduate Texts in Mathematics \textbf{218}, Springer, 2013,
\href{https://doi.org/10.1007/978-1-4419-9982-5}{doi:10.1007/978-1-4419-9982-5}.
\end{thebibliography}
\end{document}
